\section{Data and data contracts}\label{sec:data}
Every numerical input in the evaluated product is synthetic or adapted from an
identified open synthetic model. There are no private building measurements,
customer records or operational performance labels. The fifteen-case catalog
uses a fixed authored seed and separates scalar semantics, hard-row integrity,
sparse structure, predictive control and energy policy. The registered original
and grouped studies introduce their own prospective instance grids and seed
splits; neither is inferred from the teaching catalog. The code and model
attribution retain their separate upstream terms, while explanatory material
and the archived scientific datasets carry their stated distribution licences.

The C13 building source follows the published two-zone model from OSQPSlack,
with a 20~kWh battery and 96 quarter-hour intervals. It preserves mixed quadratic,
grouped, elastic and Huber prices. C14 and C15 reuse the physical structure with
explicitly changed separable side-price policies, so their optima are not
reproductions of the mixed objective. The grouped study's hourly thermal inputs
use a separate synthetic RC forecast model with hard dynamics, initial conditions
and heater-energy limits. A forecast input is not a measured trajectory or a
closed-loop controller experiment. These distinctions constrain every comparison
and figure interpretation.

The eighteen versioned boundaries below specify how those scientific objects
move between producers, numerical tools and readers. Tables are generated from
the actual JSON schemas; the accompanying prose explains the physical semantics
and cross-field checks that a type validator cannot prove. Unknown fields,
malformed dimensions and invalid states are refused according to the respective
schema and consumer. Missing data remain missing rather than becoming an empty
successful result. A schema example establishes shape, not deployed availability.
Physical units are explicit row/decision metadata where they apply; identifiers,
hashes, enum states and counts have no physical dimension.

\subsection{Deployment case inventory}\label{sec:contract-case-index}
The case-index contract schema 3.1.0 is produced by the manifest component and
consumed by the static artifact loader and optional file-serving API. Its schema,
engine version, case count, default case, case list and file inventory are shown
in Table~\ref{tab:contract-case-index}. Each listed artifact has an identity that
the reader checks before using its contents. Counts are nonnegative integers;
paths and hashes are identifiers rather than physical quantities. A default
selection must refer to a listed case, and a file's checksum must match its
actual bytes. Schema validity alone cannot establish either condition. Absence
of this inventory stops the site build because a reader cannot safely discover
its scientific evidence. An older pipeline version is stale and is displayed
as such. The index never proves that every possible solver request has run;
it names the finite exported distribution. Replacing the source with a fixture
inventory exercises the same consumer boundary without pretending the fixture
is the complete production case population.
\begin{table}[htbp]\centering\footnotesize
\caption{Schema-derived top-level fields, dimensions and literal structural restrictions for case-index schema 3.1.0. Complete nested numerical definitions are expanded in the contract appendix; physical maps remain those of the original model.}\label{tab:contract-case-index}
\begin{tabularx}{\textwidth}{@{}>{\raggedright\arraybackslash}p{.22\textwidth}>{\raggedright\arraybackslash}p{.12\textwidth}>{\raggedright\arraybackslash}X>{\raggedright\arraybackslash}X@{}}\toprule
Field & Type & Constraint & Unit / role\\\midrule
schema & string & const=\allowbreak{}softqp.\allowbreak{}artifact-\allowbreak{}index/\allowbreak{}v1 & structural metadata /\allowbreak{} nested units\\
engine\_\allowbreak{}version & string & See defined object /\allowbreak{} value type & structural metadata /\allowbreak{} nested units\\
n\_\allowbreak{}cases & integer & minimum=\allowbreak{}1 & count /\allowbreak{} integer identifier\\
default\_\allowbreak{}case & string & pattern=\allowbreak{}\textasciicircum{}C[0-9]\{2\}\$ & structural metadata /\allowbreak{} nested units\\
cases & array & minItems=\allowbreak{}1 & structural metadata /\allowbreak{} nested units\\
files & array & uniqueItems=\allowbreak{}True & structural metadata /\allowbreak{} nested units\\
\bottomrule\end{tabularx}
\end{table}


\subsection{Individual deployment manifest}\label{sec:contract-case-manifest}
The case-manifest contract schema 5.0.0 connects one scientific case to the
shared deployment shell. The manifest producer assigns the case identifier,
bilingual title/category, engine identity, product version, seed, source,
artifact descriptors, lane and acceptance information. The consumer uses that
envelope to select available evidence without embedding numerical assumptions
in navigation code. Table~\ref{tab:contract-case-manifest} lists the actual
fields. A seed is a reproducibility identifier; it is not a physical measurement.
Each artifact retains its own unit-bearing scientific payload. If an indexed
manifest is absent, readiness is partial and the corresponding case is unavailable,
while other cases remain accessible. Unknown cases receive a missing-file response.
A stale engine/product identity cannot be hidden by preserving an old checksum.
The producer and consumer tests also exercise source substitution, preventing a
handwritten browser fixture from becoming a special unvalidated path. This
contract describes presentation availability, not a certificate of every
mathematical optimum contained in the referenced payload.
\begin{table}[htbp]\centering\footnotesize
\caption{Schema-derived top-level fields, dimensions and literal structural restrictions for case-manifest schema 5.0.0. Complete nested numerical definitions are expanded in the contract appendix; physical maps remain those of the original model.}\label{tab:contract-case-manifest}
\begin{tabularx}{\textwidth}{@{}>{\raggedright\arraybackslash}p{.22\textwidth}>{\raggedright\arraybackslash}p{.12\textwidth}>{\raggedright\arraybackslash}X>{\raggedright\arraybackslash}X@{}}\toprule
Field & Type & Constraint & Unit / role\\\midrule
schema & string & const=\allowbreak{}softqp.\allowbreak{}manifest/\allowbreak{}v1 & structural metadata /\allowbreak{} nested units\\
case\_\allowbreak{}id & string & pattern=\allowbreak{}\textasciicircum{}C[0-9]\{2\}\$ & structural metadata /\allowbreak{} nested units\\
title & object & additionalProperties=\allowbreak{}False & structural metadata /\allowbreak{} nested units\\
category & object & additionalProperties=\allowbreak{}False & structural metadata /\allowbreak{} nested units\\
engine & object & additionalProperties=\allowbreak{}False & structural metadata /\allowbreak{} nested units\\
product\_\allowbreak{}version & string & See defined object /\allowbreak{} value type & structural metadata /\allowbreak{} nested units\\
seed & integer & See defined object /\allowbreak{} value type & count /\allowbreak{} integer identifier\\
source & string & See defined object /\allowbreak{} value type & structural metadata /\allowbreak{} nested units\\
artifact & object & additionalProperties=\allowbreak{}False & structural metadata /\allowbreak{} nested units\\
lane & string & const=\allowbreak{}precompute & structural metadata /\allowbreak{} nested units\\
gate & object & additionalProperties=\allowbreak{}False & structural metadata /\allowbreak{} nested units\\
\bottomrule\end{tabularx}
\end{table}


\subsection{Readiness and catalog identity}\label{sec:contract-readiness}
The readiness contract schema 2.0.1 is computed from the actual contract catalog,
product identity and distributed inventory. Its schema tag, product version,
catalog digest and per-contract states are shown in Table~\ref{tab:contract-readiness}.
It is consumed by the shared reader and optional API to distinguish available,
provider, missing, partial and stale evidence within the shared state vocabulary.
These are availability states with explanatory reasons, not accuracy labels.
The producer does not infer a trained model from an empty model directory, and
the readiness document excludes itself from recursively proving its own presence.
An absent file prevents a complete build; a runtime load failure is reported
as unknown readiness while other views retain their own evidence. The catalog
digest prevents a changed boundary inventory from being paired silently with an
old readiness result. Four grouped-study boundaries remain unavailable to the
browser because no validated replay adapter is installed. Their valid schema
examples and complete external study data do not change that missing state.
\begin{table}[htbp]\centering\footnotesize
\caption{Schema-derived top-level fields, dimensions and literal structural restrictions for readiness schema 2.0.1. Complete nested numerical definitions are expanded in the contract appendix; physical maps remain those of the original model.}\label{tab:contract-readiness}
\begin{tabularx}{\textwidth}{@{}>{\raggedright\arraybackslash}p{.22\textwidth}>{\raggedright\arraybackslash}p{.12\textwidth}>{\raggedright\arraybackslash}X>{\raggedright\arraybackslash}X@{}}\toprule
Field & Type & Constraint & Unit / role\\\midrule
schema & string & const=\allowbreak{}softqp.\allowbreak{}readiness/\allowbreak{}v2 & structural metadata /\allowbreak{} nested units\\
product\_\allowbreak{}version & string & pattern=\allowbreak{}\textasciicircum{}\textbackslash{}d+\textbackslash{}.\allowbreak{}\textbackslash{}d\{2\}\textbackslash{}.\allowbreak{}\textbackslash{}d\{3\}\$ & structural metadata /\allowbreak{} nested units\\
catalog\_\allowbreak{}sha256 & string & pattern=\allowbreak{}\textasciicircum{}[0-9a-\allowbreak{}f]\{64\}\$ & structural metadata /\allowbreak{} nested units\\
contracts & array & minItems=\allowbreak{}1 & structural metadata /\allowbreak{} nested units\\
\bottomrule\end{tabularx}
\end{table}


\subsection{Dormant file-serving API}\label{sec:contract-api}
The API contract schema 2.1.0 is an OpenAPI boundary owned by the optional app
service. In the current static product this service is dormant. Its unavailable
body has a detail field, while the full OpenAPI description also specifies
case, artifact and readiness routes and response states. Table~\ref{tab:contract-api}
describes the unavailable body rather than inventing a live inference response.
The producer serves precomputed files when explicitly activated; it never
executes an optimizer as a consequence of an HTTP request. The static reader
uses committed artifacts directly and therefore does not depend on a resident
API. A missing routed document returns the documented unavailable response
with its reason, and an unknown case or artifact is missing. An enabled
service's own readiness is distinct from the static distribution's explicit
dormant-API entry. This distinction avoids interpreting an example response,
a schema-valid object or a local development server as evidence of a deployed
numerical computation service.
\begin{table}[htbp]\centering\footnotesize
\caption{Schema-derived top-level fields, dimensions and literal structural restrictions for api schema 2.1.0. Complete nested numerical definitions are expanded in the contract appendix; physical maps remain those of the original model.}\label{tab:contract-api}
\begin{tabularx}{\textwidth}{@{}>{\raggedright\arraybackslash}p{.22\textwidth}>{\raggedright\arraybackslash}p{.12\textwidth}>{\raggedright\arraybackslash}X>{\raggedright\arraybackslash}X@{}}\toprule
Field & Type & Constraint & Unit / role\\\midrule
detail & object & additionalProperties=\allowbreak{}False & structural metadata /\allowbreak{} nested units\\
\bottomrule\end{tabularx}
\end{table}


\subsection{Scientific case inventory}\label{sec:contract-softqp-index}
The softqp-index contract schema 1.1.0 is produced by the scientific catalog
runner and consumed by the workbench's case-data loader. It carries its schema
version, overall state, default case, case descriptions, environment and
limitations; Table~\ref{tab:contract-softqp-index} gives the field inventory.
It names actual scientific records separately from shell deployment manifests.
The case count and identifiers describe a finite synthetic catalog, not a
sample of field deployments. Environment strings identify how the exported
objects were computed; they do not make an arbitrary user's local runtime
equivalent. A missing inventory leaves the scientific workbench unavailable
with a path-specific explanation. Development/provider evidence is identified
as such. Consumers compare listed paths and byte identities before interpreting
case records, and source-swap tests exercise the same code using a real
alternative fixture tree. This boundary prevents a visually populated case
selector from being confused with loaded, current and scientifically compatible
data for the selected method.
\begin{table}[htbp]\centering\footnotesize
\caption{Schema-derived top-level fields, dimensions and literal structural restrictions for softqp-index schema 1.1.0. Complete nested numerical definitions are expanded in the contract appendix; physical maps remain those of the original model.}\label{tab:contract-softqp-index}
\begin{tabularx}{\textwidth}{@{}>{\raggedright\arraybackslash}p{.22\textwidth}>{\raggedright\arraybackslash}p{.12\textwidth}>{\raggedright\arraybackslash}X>{\raggedright\arraybackslash}X@{}}\toprule
Field & Type & Constraint & Unit / role\\\midrule
schema\_\allowbreak{}version & string & const=\allowbreak{}softqp.\allowbreak{}index/\allowbreak{}v1 & structural metadata /\allowbreak{} nested units\\
state & string & const=\allowbreak{}development & structural metadata /\allowbreak{} nested units\\
default\_\allowbreak{}case & union /\allowbreak{} composition & See defined object /\allowbreak{} value type & structural metadata /\allowbreak{} nested units\\
cases & array & See defined object /\allowbreak{} value type & structural metadata /\allowbreak{} nested units\\
environment & object & additionalProperties=\allowbreak{}False & structural metadata /\allowbreak{} nested units\\
limitations & array & See defined object /\allowbreak{} value type & structural metadata /\allowbreak{} nested units\\
\bottomrule\end{tabularx}
\end{table}


\subsection{Original-space case result}\label{sec:contract-softqp-case-result}
The softqp-case-result contract schema 2.3.0 contains a case identifier, bilingual
descriptive fields, seed/source, numerical settings, method outcomes, original
row metadata and, where supplied, the sparse problem. The scientific runner
produces it; the workbench validates it and computes its physical displays.
Table~\ref{tab:contract-softqp-case-result} lists its fields. Each run retains
state and reason, original decision/dual vectors, dimensions, termination and
objective/violation metrics when available. Row units attach to original
departures, not auxiliary lifted coordinates. Method incompatibility and
numerical failure do not satisfy the available-result branch of the schema.
Cross-field consumers check vector lengths, original objectives, hard policy
and lower/upper violation consistency. A missing case leaves that case
unavailable while other records remain usable. A zero numerical hard error is
reported only when the selected unit actually contains hard rows; absence of
such rows is unavailable, not a fabricated zero. This preserves both scientific
semantics and unit-aware interpretation.
\begin{table}[htbp]\centering\footnotesize
\caption{Schema-derived top-level fields, dimensions and literal structural restrictions for softqp-case-result schema 2.3.0. Complete nested numerical definitions are expanded in the contract appendix; physical maps remain those of the original model.}\label{tab:contract-softqp-case-result}
\begin{tabularx}{\textwidth}{@{}>{\raggedright\arraybackslash}p{.22\textwidth}>{\raggedright\arraybackslash}p{.12\textwidth}>{\raggedright\arraybackslash}X>{\raggedright\arraybackslash}X@{}}\toprule
Field & Type & Constraint & Unit / role\\\midrule
schema\_\allowbreak{}version & string & const=\allowbreak{}softqp.\allowbreak{}case-\allowbreak{}result/\allowbreak{}v2 & structural metadata /\allowbreak{} nested units\\
case\_\allowbreak{}id & string & pattern=\allowbreak{}\textasciicircum{}C[0-9]\{2\}\$ & structural metadata /\allowbreak{} nested units\\
title & object & additionalProperties=\allowbreak{}False & structural metadata /\allowbreak{} nested units\\
category & object & additionalProperties=\allowbreak{}False & structural metadata /\allowbreak{} nested units\\
seed & integer & See defined object /\allowbreak{} value type & count /\allowbreak{} integer identifier\\
source & string & minLength=\allowbreak{}1 & structural metadata /\allowbreak{} nested units\\
settings & object & additionalProperties=\allowbreak{}False & structural metadata /\allowbreak{} nested units\\
runs & array & minItems=\allowbreak{}1 & structural metadata /\allowbreak{} nested units\\
rows & array & See defined object /\allowbreak{} value type & count /\allowbreak{} integer identifier\\
problem & object & additionalProperties=\allowbreak{}False & structural metadata /\allowbreak{} nested units\\
\bottomrule\end{tabularx}
\end{table}


\subsection{Sparse scientific problem specification}\label{sec:contract-softqp-case-spec}
The softqp-case-spec contract schema 1.0.0 is the original-space input exchanged
by local request execution. Its matrices $P,A$, vector $q$, lower/upper bounds
and penalty object appear in Table~\ref{tab:contract-softqp-case-spec}. Sparse
matrices use explicit compressed-column arrays: data, row indices, column
pointers and shape. Pointer monotonicity, index range, finite coefficients and
dimension agreement are numerical cross-field requirements beyond a list type.
Bounds may explicitly omit a side; missing endpoints are not replaced by
arbitrary large finite constants. The penalty object preserves hard/soft
policy and side coefficients. The producer is the local I/O component, and
the native executor and browser import validator are its consumers. A complete
matching specification is required before local execution can be accepted;
its absence does not disable canonical replay. Physical units remain supplied
by the original case's row/decision context, so a structurally valid matrix
cannot independently redefine metres as centimetres or change a hard physical
row into a finite-cost preference.
\begin{table}[htbp]\centering\footnotesize
\caption{Schema-derived top-level fields, dimensions and literal structural restrictions for softqp-case-spec schema 1.0.0. Complete nested numerical definitions are expanded in the contract appendix; physical maps remain those of the original model.}\label{tab:contract-softqp-case-spec}
\begin{tabularx}{\textwidth}{@{}>{\raggedright\arraybackslash}p{.22\textwidth}>{\raggedright\arraybackslash}p{.12\textwidth}>{\raggedright\arraybackslash}X>{\raggedright\arraybackslash}X@{}}\toprule
Field & Type & Constraint & Unit / role\\\midrule
P & object & additionalProperties=\allowbreak{}False & structural metadata /\allowbreak{} nested units\\
q & array & maxItems=\allowbreak{}20000 & cost /\allowbreak{} decision unit\\
A & object & additionalProperties=\allowbreak{}False & structural metadata /\allowbreak{} nested units\\
lower & array & maxItems=\allowbreak{}20000 & original row unit\\
upper & array & maxItems=\allowbreak{}20000 & original row unit\\
penalty & object & additionalProperties=\allowbreak{}False & structural metadata /\allowbreak{} nested units\\
\bottomrule\end{tabularx}
\end{table}


\subsection{Bounded local cost-edit request}\label{sec:contract-softqp-run-request}
The softqp-run-request contract schema 1.1.0 carries the case identity, immutable
base digest, schema version and three edit factors; its fields are in
Table~\ref{tab:contract-softqp-run-request}. The browser creates the request,
and the local scientific request runner consumes it. Lower and upper factors
range from 0.01 to 100, while the additional quadratic factor ranges from 0 to 100.
They are dimensionless multipliers on existing side prices, not replacement
physical units or arbitrary sparse matrices. The base digest binds the edit
to the exact canonical model. Cross-field validation preserves matrices,
linear term, bounds and hard-row coefficients; grouped/Huber models outside
the separable editor are refused. If no matching complete request exists,
the local path is unavailable while stored replay remains valid. A changed
digest, extra field or unsupported factor cannot be repaired by accepting an
unrelated result. This narrowly defined exchange makes a local what-if policy
change inspectable without implying a hosted solver or browser-native execution.
\begin{table}[htbp]\centering\footnotesize
\caption{Schema-derived top-level fields, dimensions and literal structural restrictions for softqp-run-request schema 1.1.0. Complete nested numerical definitions are expanded in the contract appendix; physical maps remain those of the original model.}\label{tab:contract-softqp-run-request}
\begin{tabularx}{\textwidth}{@{}>{\raggedright\arraybackslash}p{.22\textwidth}>{\raggedright\arraybackslash}p{.12\textwidth}>{\raggedright\arraybackslash}X>{\raggedright\arraybackslash}X@{}}\toprule
Field & Type & Constraint & Unit / role\\\midrule
schema\_\allowbreak{}version & string & const=\allowbreak{}softqp.\allowbreak{}run-\allowbreak{}request/\allowbreak{}v1 & structural metadata /\allowbreak{} nested units\\
case\_\allowbreak{}id & string & pattern=\allowbreak{}\textasciicircum{}C[0-9]\{2\}\$ & structural metadata /\allowbreak{} nested units\\
base\_\allowbreak{}sha256 & string & pattern=\allowbreak{}\textasciicircum{}[0-9a-\allowbreak{}f]\{64\}\$ & structural metadata /\allowbreak{} nested units\\
lower\_\allowbreak{}multiplier & number & minimum=\allowbreak{}0.\allowbreak{}01; maximum=\allowbreak{}100 & dimensionless factor\\
upper\_\allowbreak{}multiplier & number & minimum=\allowbreak{}0.\allowbreak{}01; maximum=\allowbreak{}100 & dimensionless factor\\
quadratic\_\allowbreak{}multiplier & number & minimum=\allowbreak{}0; maximum=\allowbreak{}100 & dimensionless factor\\
\bottomrule\end{tabularx}
\end{table}


\subsection{Local native result envelope}\label{sec:contract-softqp-local-export}
The softqp-local-export contract schema 1.1.0 wraps the original request,
result, schema identity and optional runtime information. The local request
runner produces it, and the browser verifies it before changing the selected
workbench result. Table~\ref{tab:contract-softqp-local-export} gives the top-level
fields; nested request and result retain their own complete definitions.
Runtime fingerprints improve provenance but do not replace original-space
numerical checks. The consumer matches case digest and edit factors, checks
unchanged hard policy, recomputes objective components and departures, and
rejects altered imported objectives. Missing or mismatched exports leave
local execution unavailable, preserving canonical replay. Accepted imports
are labelled local-native and can be restored to the original record; changing
the selected case discards an incompatible local import. No upload is sent
to a hosted service. The exported file is therefore evidence of a particular
external numerical execution under a matching request, not proof that the
reader itself computed an optimizer or that every result from a user machine
has the historical experiment's source identity.
\begin{table}[htbp]\centering\footnotesize
\caption{Schema-derived top-level fields, dimensions and literal structural restrictions for softqp-local-export schema 1.1.0. Complete nested numerical definitions are expanded in the contract appendix; physical maps remain those of the original model.}\label{tab:contract-softqp-local-export}
\begin{tabularx}{\textwidth}{@{}>{\raggedright\arraybackslash}p{.22\textwidth}>{\raggedright\arraybackslash}p{.12\textwidth}>{\raggedright\arraybackslash}X>{\raggedright\arraybackslash}X@{}}\toprule
Field & Type & Constraint & Unit / role\\\midrule
schema\_\allowbreak{}version & string & const=\allowbreak{}softqp.\allowbreak{}local-\allowbreak{}export/\allowbreak{}v1 & structural metadata /\allowbreak{} nested units\\
request & object & additionalProperties=\allowbreak{}False & structural metadata /\allowbreak{} nested units\\
result & object & additionalProperties=\allowbreak{}False & structural metadata /\allowbreak{} nested units\\
environment & object & additionalProperties=\allowbreak{}False & structural metadata /\allowbreak{} nested units\\
\bottomrule\end{tabularx}
\end{table}


\subsection{Paired predictive-control stream}\label{sec:contract-softqp-stream}
The softqp-stream contract schema 1.0.0 is produced by the experiment exporter
and consumed by the fixed C12 stream reader. It carries case/protocol/source
identities, environment, stream specification, frames, counts and limitations,
as shown in Table~\ref{tab:contract-softqp-stream}. Frames retain physical time,
the common state, oracle action, disturbances and method/mode observations.
Position, velocity and acceleration use metres, metres per second and metres
per squared second; elapsed numerical durations use seconds. Thirty-two
frames share one independently driven trajectory, so comparing their solver
outputs does not compare separate closed-loop controllers. Warm and cold
workspaces have explicit setup/update boundaries. A missing stream is shown
as missing with its path without erasing other evidence. The stream remains
unavailable after an edited local model because its recorded policy no longer
matches that request. Full frame ordering and mode identities are checked;
overlapping predicted horizon costs cannot be summed as realized operational
costs, and incompatible generalized projection states remain visible.
\begin{table}[htbp]\centering\footnotesize
\caption{Schema-derived top-level fields, dimensions and literal structural restrictions for softqp-stream schema 1.0.0. Complete nested numerical definitions are expanded in the contract appendix; physical maps remain those of the original model.}\label{tab:contract-softqp-stream}
\begin{tabularx}{\textwidth}{@{}>{\raggedright\arraybackslash}p{.22\textwidth}>{\raggedright\arraybackslash}p{.12\textwidth}>{\raggedright\arraybackslash}X>{\raggedright\arraybackslash}X@{}}\toprule
Field & Type & Constraint & Unit / role\\\midrule
schema\_\allowbreak{}version & string & const=\allowbreak{}softqp.\allowbreak{}stream/\allowbreak{}v1 & structural metadata /\allowbreak{} nested units\\
case\_\allowbreak{}id & string & const=\allowbreak{}C12 & structural metadata /\allowbreak{} nested units\\
protocol\_\allowbreak{}id & string & const=\allowbreak{}softqp-\allowbreak{}experiments-\allowbreak{}v1 & structural metadata /\allowbreak{} nested units\\
engine\_\allowbreak{}source & string & pattern=\allowbreak{}\textasciicircum{}[0-9a-\allowbreak{}f]\{40\}\$ & structural metadata /\allowbreak{} nested units\\
environment & object & additionalProperties=\allowbreak{}False & structural metadata /\allowbreak{} nested units\\
specification & object & additionalProperties=\allowbreak{}False & structural metadata /\allowbreak{} nested units\\
frames & array & minItems=\allowbreak{}32; maxItems=\allowbreak{}32 & count /\allowbreak{} integer identifier\\
counts & object & additionalProperties=\allowbreak{}False & structural metadata /\allowbreak{} nested units\\
limitations & array & See defined object /\allowbreak{} value type & structural metadata /\allowbreak{} nested units\\
\bottomrule\end{tabularx}
\end{table}


\subsection{Original experiment summary}\label{sec:contract-softqp-campaign} % voice-ok: literal contract schema identifier or source path
The softqp-campaign contract schema 1.0.0 summarizes the original registered % voice-ok: literal contract schema identifier or source path
study. The exporter supplies the protocol, environment, source/digest maps,
observed process duration, peak worker memory, restart count, planned/completed
instance counts, cell summaries, primary endpoint and linked artifacts.
Table~\ref{tab:contract-softqp-campaign} lists these fields. Durations are in % voice-ok: literal contract schema identifier or source path
seconds and resident memory in bytes; objective discrepancies and relative
reductions are dimensionless. The browser's experimental reader consumes the
summary after byte verification. A completed count must agree with the actual
instance inventory and retained trial states. Its primary endpoint is the
stored original 168-pair analysis, not a value recomputed from whichever
descriptive plot is currently filtered. Missing evidence is explicitly
unavailable while other case data remain present. Runtime measurements describe
the specified supervised process and are not process-wide maxima inferred
from CPU activity. This contract connects numerical observations to their
registered uncertainty without treating a schema-valid positive replacement
or a newly written checksum as independent scientific acceptance.
\begin{table}[htbp]\centering\footnotesize
\caption{Schema-derived top-level fields, dimensions and literal structural restrictions for softqp-campaign schema 1.0.0. Complete nested numerical definitions are expanded in the contract appendix; physical maps remain those of the original model.}\label{tab:contract-softqp-campaign} % voice-ok: literal schema identifier or archived schema value, not manuscript narrative
\begin{tabularx}{\textwidth}{@{}>{\raggedright\arraybackslash}p{.22\textwidth}>{\raggedright\arraybackslash}p{.12\textwidth}>{\raggedright\arraybackslash}X>{\raggedright\arraybackslash}X@{}}\toprule
Field & Type & Constraint & Unit / role\\\midrule
schema\_\allowbreak{}version & string & const=\allowbreak{}softqp.\allowbreak{}campaign/\allowbreak{}v1 & structural metadata /\allowbreak{} nested units\\ % voice-ok: literal schema identifier or archived schema value, not manuscript narrative
protocol & object & additionalProperties=\allowbreak{}False & structural metadata /\allowbreak{} nested units\\
environment & object & additionalProperties=\allowbreak{}False & structural metadata /\allowbreak{} nested units\\
protocol\_\allowbreak{}sha256 & string & pattern=\allowbreak{}\textasciicircum{}[0-9a-\allowbreak{}f]\{64\}\$ & structural metadata /\allowbreak{} nested units\\
source\_\allowbreak{}sha256 & object & additionalProperties=\allowbreak{}\{'type': 'string', 'pattern': '\textasciicircum{}[0-9a-\allowbreak{}f]\{64\}\$'\} & structural metadata /\allowbreak{} nested units\\
observed\_\allowbreak{}wall\_\allowbreak{}seconds & number & minimum=\allowbreak{}0 & s\\
worker\_\allowbreak{}peak\_\allowbreak{}rss\_\allowbreak{}bytes & integer & minimum=\allowbreak{}0 & bytes\\
worker\_\allowbreak{}restarts & integer & minimum=\allowbreak{}0 & count /\allowbreak{} integer identifier\\
planned\_\allowbreak{}instances & integer & minimum=\allowbreak{}0 & count /\allowbreak{} integer identifier\\
completed\_\allowbreak{}instances & integer & minimum=\allowbreak{}0 & count /\allowbreak{} integer identifier\\
cells & array & See defined object /\allowbreak{} value type & structural metadata /\allowbreak{} nested units\\
primary\_\allowbreak{}endpoint & object & additionalProperties=\allowbreak{}False & structural metadata /\allowbreak{} nested units\\
counts & object & additionalProperties=\allowbreak{}False & structural metadata /\allowbreak{} nested units\\
artifacts & array & See defined object /\allowbreak{} value type & structural metadata /\allowbreak{} nested units\\
\bottomrule\end{tabularx}
\end{table}
 % voice-ok: literal contract schema identifier or source path

\subsection{Seeded instance inventory}\label{sec:contract-softqp-datasets}
The softqp-datasets contract schema 1.0.0 governs the original study's array
of generated-instance descriptors. Each entry identifies its split, dimension,
density, diagonal condition, penalty family, seed, generation duration,
problem digest and independent oracle state/reason/objective. Its item fields
are shown in Table~\ref{tab:contract-softqp-datasets}. This is a list contract,
not an object whose missing top-level properties imply an empty dataset.
The experiment exporter produces it; the benchmark reader and replay validator
consume it. Seeds and digests identify exact generated models. Density and
condition are dimensionless controlled inputs; generation duration is a
reported process interval distinct from solver time. A missing list makes
experiment evidence unavailable with its specific path. Duplicate or omitted
instances are inconsistent with the prospective population even if every
remaining entry validates separately. Oracle failure is retained rather than
replacing its model with an easier seed. This inventory therefore supports
denominator and compatibility checks without turning an authored synthetic
case grid into an operational sampling claim.
\begin{table}[htbp]\centering\footnotesize
\caption{Schema-derived top-level fields, dimensions and literal structural restrictions for softqp-datasets schema 1.0.0. Complete nested numerical definitions are expanded in the contract appendix; physical maps remain those of the original model.}\label{tab:contract-softqp-datasets}
\begin{tabularx}{\textwidth}{@{}>{\raggedright\arraybackslash}p{.22\textwidth}>{\raggedright\arraybackslash}p{.12\textwidth}>{\raggedright\arraybackslash}X>{\raggedright\arraybackslash}X@{}}\toprule
Field & Type & Constraint & Unit / role\\\midrule
id & string & pattern=\allowbreak{}\textasciicircum{}(development\textbar{}\allowbreak{}heldout\textbar{}\allowbreak{}robustness)-[0-9]\{4\}\$ & structural metadata /\allowbreak{} nested units\\
split & enumerated value & enum=\allowbreak{}['development', 'heldout', 'robustness'] & structural metadata /\allowbreak{} nested units\\
n & integer & minimum=\allowbreak{}0 & count /\allowbreak{} integer identifier\\
density & number & minimum=\allowbreak{}0 & dimensionless normalized diagnostic /\allowbreak{} factor\\
condition & number & minimum=\allowbreak{}0 & dimensionless normalized diagnostic /\allowbreak{} factor\\
family & enumerated value & enum=\allowbreak{}['symmetric-\allowbreak{}elastic', 'asymmetric-\allowbreak{}elastic', 'asymmetric-\allowbreak{}l1'] & structural metadata /\allowbreak{} nested units\\
seed & integer & minimum=\allowbreak{}0 & count /\allowbreak{} integer identifier\\
generation\_\allowbreak{}seconds & number & minimum=\allowbreak{}0 & s\\
problem\_\allowbreak{}sha256 & string & pattern=\allowbreak{}\textasciicircum{}[0-9a-\allowbreak{}f]\{64\}\$ & structural metadata /\allowbreak{} nested units\\
oracle\_\allowbreak{}state & enumerated value & enum=\allowbreak{}['accepted', 'oracle-\allowbreak{}failed'] & structural metadata /\allowbreak{} nested units\\
oracle\_\allowbreak{}reason & union /\allowbreak{} composition & See defined object /\allowbreak{} value type & structural metadata /\allowbreak{} nested units\\
oracle\_\allowbreak{}objective & union /\allowbreak{} composition & See defined object /\allowbreak{} value type & model cost unit\\
\bottomrule\end{tabularx}
\end{table}


\subsection{Original per-trial observation}\label{sec:contract-softqp-trial}
The softqp-trial contract schema 1.0.0 governs each JSONL observation and its
CSV projection. It retains the instance/split/family controls, method,
repetition and warm-up flag, state/reason, original metrics, discrepancy,
elapsed/update/setup durations, numerical result information, iterations,
dimensions and backend status fields. Table~\ref{tab:contract-softqp-trial}
lists the actual field inventory. The experiment exporter produces these
rows; replay and benchmark consumers preserve their identity and type.
Numerical durations are seconds, matrix sizes are counts and acceptance
metrics have their stated normalized or physical interpretation. An
incompatible observation can carry no solved objective; absence is not
encoded as a successful zero. Warm-up flags prevent initialization samples
from entering the measured primary reduction. Missing raw evidence remains
missing even if a summary CSV exists. Full JSONL retains diagnostics that
cannot be reconstructed from selected plot points. Repetition, method
ordering and reason fields are retained so a fast unsuccessful or unsupported
call cannot silently enter an accuracy-qualified timing distribution.
\begingroup\footnotesize
\begin{longtable}{@{}>{\raggedright\arraybackslash}p{.22\textwidth}>{\raggedright\arraybackslash}p{.12\textwidth}>{\raggedright\arraybackslash}p{.32\textwidth}>{\raggedright\arraybackslash}p{.25\textwidth}@{}}
\caption{Schema-derived top-level fields, dimensions and literal structural restrictions for softqp-trial schema 1.0.0. Complete nested numerical definitions are expanded in the contract appendix; physical maps remain those of the original model.}\label{tab:contract-softqp-trial}\\
\toprule Field & Type & Constraint & Unit / role\\\midrule\endfirsthead
\toprule Field & Type & Constraint & Unit / role\\\midrule\endhead
\bottomrule\endfoot
id & string & pattern=\allowbreak{}\textasciicircum{}(development\textbar{}\allowbreak{}heldout\textbar{}\allowbreak{}robustness)-[0-9]\{4\}\$ & structural metadata /\allowbreak{} nested units\\
split & enumerated value & enum=\allowbreak{}['development', 'heldout', 'robustness'] & structural metadata /\allowbreak{} nested units\\
n & integer & minimum=\allowbreak{}0 & count /\allowbreak{} integer identifier\\
density & number & minimum=\allowbreak{}0 & dimensionless normalized diagnostic /\allowbreak{} factor\\
condition & number & minimum=\allowbreak{}0 & dimensionless normalized diagnostic /\allowbreak{} factor\\
family & enumerated value & enum=\allowbreak{}['symmetric-\allowbreak{}elastic', 'asymmetric-\allowbreak{}elastic', 'asymmetric-\allowbreak{}l1'] & structural metadata /\allowbreak{} nested units\\
seed & integer & minimum=\allowbreak{}0 & count /\allowbreak{} integer identifier\\
method & enumerated value & enum=\allowbreak{}['lifted', 'generalized', 'asymmetric', 'daqp', 'flex', 'elastiqp'] & structural metadata /\allowbreak{} nested units\\
repetition & integer & minimum=\allowbreak{}-1 & count /\allowbreak{} integer identifier\\
warmup & boolean & See defined object /\allowbreak{} value type & structural metadata /\allowbreak{} nested units\\
state & enumerated value & enum=\allowbreak{}['accepted', 'incompatible', 'failed', 'accuracy-\allowbreak{}rejected', 'timeout', 'resource-\allowbreak{}limit', 'oracle-\allowbreak{}failed'] & structural metadata /\allowbreak{} nested units\\
reason & union /\allowbreak{} composition & See defined object /\allowbreak{} value type & structural metadata /\allowbreak{} nested units\\
metrics & union /\allowbreak{} composition & See defined object /\allowbreak{} value type & structural metadata /\allowbreak{} nested units\\
normalized\_\allowbreak{}gap & union /\allowbreak{} composition & See defined object /\allowbreak{} value type & dimensionless normalized diagnostic /\allowbreak{} factor\\
elapsed\_\allowbreak{}seconds & union /\allowbreak{} composition & See defined object /\allowbreak{} value type & s\\
update\_\allowbreak{}seconds & union /\allowbreak{} composition & See defined object /\allowbreak{} value type & s\\
workspace\_\allowbreak{}setup\_\allowbreak{}seconds & union /\allowbreak{} composition & See defined object /\allowbreak{} value type & s\\
result & union /\allowbreak{} composition & See defined object /\allowbreak{} value type & structural metadata /\allowbreak{} nested units\\
iterations & union /\allowbreak{} composition & See defined object /\allowbreak{} value type & count /\allowbreak{} integer identifier\\
dimensions & union /\allowbreak{} composition & See defined object /\allowbreak{} value type & structural metadata /\allowbreak{} nested units\\
native\_\allowbreak{}setup\_\allowbreak{}seconds & union /\allowbreak{} composition & See defined object /\allowbreak{} value type & s\\
native\_\allowbreak{}solve\_\allowbreak{}seconds & union /\allowbreak{} composition & See defined object /\allowbreak{} value type & s\\
status & union /\allowbreak{} composition & See defined object /\allowbreak{} value type & structural metadata /\allowbreak{} nested units\\
status\_\allowbreak{}code & union /\allowbreak{} composition & See defined object /\allowbreak{} value type & count /\allowbreak{} integer identifier\\
\end{longtable}\endgroup


\subsection{Experimental artifact inventory}\label{sec:contract-softqp-experiment-index}
The softqp-experiment-index contract schema 1.0.0 associates the protocol
identity and experiment state with the exported artifact list. The producer
is the experiment export module; the static experimental loader is its
consumer. Table~\ref{tab:contract-softqp-experiment-index} lists its schema
version, protocol identifier, artifacts and state. Each artifact is a named
file with a declared representation and checksum rather than an inferred
availability flag. Counts and byte identities are dimensionless metadata;
physical units remain in the scientific records. Missing inventory leaves
experiment evidence unavailable but does not disable independent teaching
cases. The browser verifies actual bytes before reading JSON or plotting
the selected projection. A valid inventory cannot certify the underlying
mathematical acceptance on its own: replay must also verify complete
populations, numerical states and source identities. Conversely, an existing
large external archive does not automatically become a declared browser
artifact. This boundary keeps the small public reader's exact distributed
files separate from external raw study storage and prevents accidental
claims that an unimplemented adapter has loaded an archived computation.
\begin{table}[htbp]\centering\footnotesize
\caption{Schema-derived top-level fields, dimensions and literal structural restrictions for softqp-experiment-index schema 1.0.0. Complete nested numerical definitions are expanded in the contract appendix; physical maps remain those of the original model.}\label{tab:contract-softqp-experiment-index}
\begin{tabularx}{\textwidth}{@{}>{\raggedright\arraybackslash}p{.22\textwidth}>{\raggedright\arraybackslash}p{.12\textwidth}>{\raggedright\arraybackslash}X>{\raggedright\arraybackslash}X@{}}\toprule
Field & Type & Constraint & Unit / role\\\midrule
schema\_\allowbreak{}version & string & const=\allowbreak{}softqp.\allowbreak{}experiments/\allowbreak{}v1 & structural metadata /\allowbreak{} nested units\\
protocol\_\allowbreak{}id & string & const=\allowbreak{}softqp-\allowbreak{}experiments-\allowbreak{}v1 & structural metadata /\allowbreak{} nested units\\
artifacts & array & See defined object /\allowbreak{} value type & structural metadata /\allowbreak{} nested units\\
state & enumerated value & enum=\allowbreak{}['development', 'measured'] & structural metadata /\allowbreak{} nested units\\
\bottomrule\end{tabularx}
\end{table}


\subsection{Directional repeated-forecast input batch}\label{sec:contract-softqp-coupled-batch}
The softqp-coupled-batch contract schema 1.0.0 is an offline research input
boundary produced and consumed by the coupled contract/worker components.
Its schema version, kind, description, original problem, vector updates and
unit arrays appear in Table~\ref{tab:contract-softqp-coupled-batch}. The
problem retains sparse matrices, hard policy, disjoint groups, positive
directional scales and original physical metadata. Updates change only
the admitted repeated-operation vectors while preserving the fixed structural
objects. Units distinguish each decision, interval row and group peak, so
an inverse-row-unit directional scale can be checked against the chosen
dimensionless peak. The batch is hashed before worker execution and
cross-referenced by its observations. A valid batch example does not
establish an available browser replay path: that consumer is absent and
readiness remains missing for this boundary. Invalid group overlap,
nonpositive scales, changed hard policy or mismatched update dimensions
are refused rather than coerced into a superficially feasible model. The
complete external archive contains the actual synthetic batches separately
from these shape examples.
\begin{table}[htbp]\centering\footnotesize
\caption{Schema-derived top-level fields, dimensions and literal structural restrictions for softqp-coupled-batch schema 1.0.0. Complete nested numerical definitions are expanded in the contract appendix; physical maps remain those of the original model.}\label{tab:contract-softqp-coupled-batch}
\begin{tabularx}{\textwidth}{@{}>{\raggedright\arraybackslash}p{.22\textwidth}>{\raggedright\arraybackslash}p{.12\textwidth}>{\raggedright\arraybackslash}X>{\raggedright\arraybackslash}X@{}}\toprule
Field & Type & Constraint & Unit / role\\\midrule
schema\_\allowbreak{}version & string & const=\allowbreak{}1.\allowbreak{}0.\allowbreak{}0 & structural metadata /\allowbreak{} nested units\\
kind & string & const=\allowbreak{}directional-\allowbreak{}coupled-\allowbreak{}peak-\allowbreak{}batch & structural metadata /\allowbreak{} nested units\\
description & string & minLength=\allowbreak{}1 & structural metadata /\allowbreak{} nested units\\
problem & object & additionalProperties=\allowbreak{}False & structural metadata /\allowbreak{} nested units\\
updates & array & minItems=\allowbreak{}1; maxItems=\allowbreak{}512 & structural metadata /\allowbreak{} nested units\\
units & object & additionalProperties=\allowbreak{}\{'type': 'string', 'minLength': 1\} & structural metadata /\allowbreak{} nested units\\
\bottomrule\end{tabularx}
\end{table}


\subsection{Directional method observation}\label{sec:contract-softqp-coupled-observation}
The softqp-coupled-observation contract schema 1.0.0 represents one oracle or
method operation in the grouped study. It retains role, instance/split/seed,
family/curvature/quadratic controls, method, repetition, warm-up, state/reason,
complete elapsed duration, input batch digest, compatibility status, frame
records, backend results, objectives and structural dimensions. The fields
are shown in Table~\ref{tab:contract-softqp-coupled-observation}. The supervised
worker produces the observation; the replay validator checks every returned
frame in original coordinates. An elapsed duration measures preparation
plus all 32 updates, solves, recovery and required acceptance calculations,
not one forecast frame. An inaccurate candidate, dependent oracle failure
and timeout remain distinct states. Missing primal/dual vectors are not
fabricated from an objective value. No browser adapter is installed for
these observations, so schema examples do not make the lane available.
The batch digest, fixed method settings and complete frame count prevent
transplanting an accepted solution from another instance or counting only
successful frames from an otherwise failed repeated operation.
\begingroup\footnotesize
\begin{longtable}{@{}>{\raggedright\arraybackslash}p{.22\textwidth}>{\raggedright\arraybackslash}p{.12\textwidth}>{\raggedright\arraybackslash}p{.32\textwidth}>{\raggedright\arraybackslash}p{.25\textwidth}@{}}
\caption{Schema-derived top-level fields, dimensions and literal structural restrictions for softqp-coupled-observation schema 1.0.0. Complete nested numerical definitions are expanded in the contract appendix; physical maps remain those of the original model.}\label{tab:contract-softqp-coupled-observation}\\
\toprule Field & Type & Constraint & Unit / role\\\midrule\endfirsthead
\toprule Field & Type & Constraint & Unit / role\\\midrule\endhead
\bottomrule\endfoot
schema\_\allowbreak{}version & string & const=\allowbreak{}1.\allowbreak{}0.\allowbreak{}0 & structural metadata /\allowbreak{} nested units\\
kind & string & const=\allowbreak{}directional-\allowbreak{}coupled-\allowbreak{}observation & structural metadata /\allowbreak{} nested units\\
role & enumerated value & enum=\allowbreak{}['oracle', 'method'] & structural metadata /\allowbreak{} nested units\\
id & string & pattern=\allowbreak{}\textasciicircum{}(development\textbar{}\allowbreak{}heldout\textbar{}\allowbreak{}robustness)-[0-9]\{4\}\$ & structural metadata /\allowbreak{} nested units\\
split & enumerated value & enum=\allowbreak{}['development', 'heldout', 'robustness'] & structural metadata /\allowbreak{} nested units\\
seed & integer & minimum=\allowbreak{}0 & count /\allowbreak{} integer identifier\\
family & enumerated value & enum=\allowbreak{}['sparse', 'thermal'] & structural metadata /\allowbreak{} nested units\\
curvature & enumerated value & enum=\allowbreak{}['positive-\allowbreak{}definite', 'semidefinite', 'zero-\allowbreak{}hessian'] & structural metadata /\allowbreak{} nested units\\
quadratic & number & minimum=\allowbreak{}0 & cost /\allowbreak{} peak unit squared\\
method & enumerated value & enum=\allowbreak{}['\_\allowbreak{}\_\allowbreak{}oracle\_\allowbreak{}\_\allowbreak{}', 'native', 'epigraph-\allowbreak{}default', 'epigraph-\allowbreak{}fixed', 'row-\allowbreak{}split-\allowbreak{}default', 'row-\allowbreak{}split-\allowbreak{}fixed', 'daqp-\allowbreak{}direct', 'daqp-\allowbreak{}prox-\allowbreak{}auto', 'daqp-\allowbreak{}prox-\allowbreak{}forced', 'daqp-\allowbreak{}lp-\allowbreak{}model'] & structural metadata /\allowbreak{} nested units\\
repetition & integer /\allowbreak{} null & minimum=\allowbreak{}-1; maximum=\allowbreak{}2 & count /\allowbreak{} integer identifier\\
warmup & boolean & See defined object /\allowbreak{} value type & structural metadata /\allowbreak{} nested units\\
state & enumerated value & enum=\allowbreak{}['accepted', 'failed', 'accuracy-\allowbreak{}rejected', 'incompatible', 'oracle-\allowbreak{}failed', 'resource-\allowbreak{}limit', 'timeout'] & structural metadata /\allowbreak{} nested units\\
reason & string /\allowbreak{} null & See defined object /\allowbreak{} value type & structural metadata /\allowbreak{} nested units\\
elapsed\_\allowbreak{}seconds & number /\allowbreak{} null & exclusiveMinimum=\allowbreak{}0 & s\\
batch\_\allowbreak{}sha256 & string & pattern=\allowbreak{}\textasciicircum{}[0-9a-\allowbreak{}f]\{64\}\$ & structural metadata /\allowbreak{} nested units\\
predeclared\_\allowbreak{}incompatible & boolean & See defined object /\allowbreak{} value type & structural metadata /\allowbreak{} nested units\\
frames & array & maxItems=\allowbreak{}32 & count /\allowbreak{} integer identifier\\
backend\_\allowbreak{}results & array & maxItems=\allowbreak{}32 & structural metadata /\allowbreak{} nested units\\
objectives & array /\allowbreak{} null & maxItems=\allowbreak{}32 & model cost unit\\
n & integer & minimum=\allowbreak{}1 & count /\allowbreak{} integer identifier\\
group\_\allowbreak{}size & integer & minimum=\allowbreak{}1 & count /\allowbreak{} integer identifier\\
condition & number & minimum=\allowbreak{}1 & dimensionless normalized diagnostic /\allowbreak{} factor\\
horizon & integer & minimum=\allowbreak{}1 & count /\allowbreak{} integer identifier\\
zones & integer & minimum=\allowbreak{}1 & count /\allowbreak{} integer identifier\\
\end{longtable}\endgroup


\subsection{Directional split summary}\label{sec:contract-softqp-coupled-campaign} % voice-ok: literal contract schema identifier or source path
The softqp-coupled-campaign contract schema 1.0.0 summarizes each registered % voice-ok: literal contract schema identifier or source path
grouped-study split. Its environment/source maps, full instance inventory,
dataset paths, measured state counts, primary fields, peak worker memory,
restart count and raw/CSV hashes are listed in
Table~\ref{tab:contract-softqp-coupled-campaign}. The offline exporter produces % voice-ok: literal contract schema identifier or source path
the summary, and independent replay consumes it. Development, held-out and
robustness populations remain separate; they cannot be pooled to replace
the 144-member primary endpoint. Elapsed quantities in nested observations
use seconds, resident memory uses bytes and structural sizes are counts.
Split names and digests identify populations rather than statistical
weights selected after timing. A browser consumer is explicitly absent,
and readiness remains missing. The external validation checks all measured
states, inventories and artifact hashes before accepting the summary.
Rewriting a checksum alongside a changed primary number therefore does
not recreate the independent evidence authority. This contract supports
archival reproducibility and denominator accounting without declaring a
hosted research execution service or a successful computational hypothesis.
\begin{table}[htbp]\centering\footnotesize
\caption{Schema-derived top-level fields, dimensions and literal structural restrictions for softqp-coupled-campaign schema 1.0.0. Complete nested numerical definitions are expanded in the contract appendix; physical maps remain those of the original model.}\label{tab:contract-softqp-coupled-campaign} % voice-ok: literal schema identifier or archived schema value, not manuscript narrative
\begin{tabularx}{\textwidth}{@{}>{\raggedright\arraybackslash}p{.22\textwidth}>{\raggedright\arraybackslash}p{.12\textwidth}>{\raggedright\arraybackslash}X>{\raggedright\arraybackslash}X@{}}\toprule
Field & Type & Constraint & Unit / role\\\midrule
schema\_\allowbreak{}version & string & const=\allowbreak{}1.\allowbreak{}0.\allowbreak{}0 & structural metadata /\allowbreak{} nested units\\
kind & string & const=\allowbreak{}directional-\allowbreak{}coupled-\allowbreak{}campaign & structural metadata /\allowbreak{} nested units\\ % voice-ok: literal schema identifier or archived schema value, not manuscript narrative
protocol\_\allowbreak{}id & string & const=\allowbreak{}softqp-\allowbreak{}directional-\allowbreak{}coupled-\allowbreak{}v1 & structural metadata /\allowbreak{} nested units\\
split & enumerated value & enum=\allowbreak{}['development', 'heldout', 'robustness'] & structural metadata /\allowbreak{} nested units\\
environment & object & See defined object /\allowbreak{} value type & structural metadata /\allowbreak{} nested units\\
source\_\allowbreak{}sha256 & object & additionalProperties=\allowbreak{}\{'\$ref': '\#/\allowbreak{}\$defs/\allowbreak{}hash'\} & structural metadata /\allowbreak{} nested units\\
inventory & array & See defined object /\allowbreak{} value type & structural metadata /\allowbreak{} nested units\\
datasets & array & See defined object /\allowbreak{} value type & structural metadata /\allowbreak{} nested units\\
counts & object & additionalProperties=\allowbreak{}\{'type': 'integer', 'minimum': 0\} & structural metadata /\allowbreak{} nested units\\
primary & object & See defined object /\allowbreak{} value type & structural metadata /\allowbreak{} nested units\\
peak\_\allowbreak{}worker\_\allowbreak{}rss & integer & minimum=\allowbreak{}0 & bytes\\
worker\_\allowbreak{}restarts & integer & minimum=\allowbreak{}0 & count /\allowbreak{} integer identifier\\
raw\_\allowbreak{}sha256 & string & pattern=\allowbreak{}\textasciicircum{}[0-9a-\allowbreak{}f]\{64\}\$ & structural metadata /\allowbreak{} nested units\\
csv\_\allowbreak{}sha256 & string & pattern=\allowbreak{}\textasciicircum{}[0-9a-\allowbreak{}f]\{64\}\$ & structural metadata /\allowbreak{} nested units\\
\bottomrule\end{tabularx}
\end{table}
 % voice-ok: literal contract schema identifier or source path

\subsection{Isolated directional teaching exchange}\label{sec:contract-softqp-coupled-lesson}
The softqp-coupled-lesson contract schema 1.0.0 connects a small physical
teaching input to actual native, compact-epigraph and directional row-split
answers. It carries schema/SDK identity, worker fingerprint, input digest,
the original input and method records; Table~\ref{tab:contract-softqp-coupled-lesson}
lists these fields. The separately configured direct-CPython worker produces
the result, and the teaching consumer verifies it before interpreting
objectives, physical decisions, group peaks, hard errors and original KKT
diagnostics. Inputs specify linear/quadratic group prices, row scaling
and a directional policy scale. Row conversion is distinct from changing
that policy. The source fingerprint and installed public SDK manifest
prevent an unrelated local module from silently replacing the tested
implementation. The exchange has bounded process time, output size and
cleanup behavior. It remains unavailable to the browser because the
teaching consumer is offline. Missing or mismatched results are refused,
and a local shape example never serves as proof of a deployed numerical
lane. This small illustration is separate from the full registered timing
population and contributes no additional inference to that endpoint.
\begin{table}[htbp]\centering\footnotesize
\caption{Schema-derived top-level fields, dimensions and literal structural restrictions for softqp-coupled-lesson schema 1.0.0. Complete nested numerical definitions are expanded in the contract appendix; physical maps remain those of the original model.}\label{tab:contract-softqp-coupled-lesson}
\begin{tabularx}{\textwidth}{@{}>{\raggedright\arraybackslash}p{.22\textwidth}>{\raggedright\arraybackslash}p{.12\textwidth}>{\raggedright\arraybackslash}X>{\raggedright\arraybackslash}X@{}}\toprule
Field & Type & Constraint & Unit / role\\\midrule
schema\_\allowbreak{}version & string & const=\allowbreak{}softqp.\allowbreak{}coupled-\allowbreak{}lesson/\allowbreak{}v1 & structural metadata /\allowbreak{} nested units\\
sdk\_\allowbreak{}version & string & const=\allowbreak{}0.\allowbreak{}2.\allowbreak{}0 & structural metadata /\allowbreak{} nested units\\
worker\_\allowbreak{}sha256 & string & pattern=\allowbreak{}\textasciicircum{}[0-9a-\allowbreak{}f]\{64\}\$ & structural metadata /\allowbreak{} nested units\\
input\_\allowbreak{}sha256 & string & pattern=\allowbreak{}\textasciicircum{}[0-9a-\allowbreak{}f]\{64\}\$ & structural metadata /\allowbreak{} nested units\\
input & object & additionalProperties=\allowbreak{}False & structural metadata /\allowbreak{} nested units\\
records & array & minItems=\allowbreak{}3; maxItems=\allowbreak{}3 & structural metadata /\allowbreak{} nested units\\
\bottomrule\end{tabularx}
\end{table}

